\subsection{Deep reinforcement learning}
\label{sec:deep}

Table~\ref{tab:deep} repeats the standard-Kyle experiment with residual memory
for Double DQN and PPO ($4\times10^5$ periods, 50 markets each; the deep agents
are far more expensive per step). Neither produces a supracompetitive outcome:
PPO ends at the competitive level ($\Delta=\PpoDeltaInt$) and DQN slightly
\emph{over}-trades ($\Delta=\DqnDeltaInt$), with myopic variants close by
($\PpoMyopicDeltaInt$ and $\DqnMyopicDeltaInt$). No variant punishes deviations.
This is what the mechanism of Section~\ref{sec:mechanism} predicts. Pruning
requires estimates that stay frozen for orders that are rarely chosen. PPO
learns from every sampled order through the policy gradient, with an entropy
bonus that keeps sampling broad, and DQN's network generalises across
neighbouring orders and reuses past experience through replay, so neither
leaves individual orders stranded with pessimistic estimates. We conjecture
that DQN's residual over-estimation from the max operator
\citep{thrun1993,vanhasselt2016}, which favours orders with noisier payoffs,
explains its slight over-trading. The deep-RL horizon is short and the
hyperparameters standard, so we view these results as evidence that the
supracompetitive outcome is specific to tabular, constant-step value learning,
not as a definitive statement about deep RL in markets
\citep[cf.][]{deng2024}.

\begin{table}[t]
  \centering
  \small
  \resizebox{0.9\linewidth}{!}{\begin{tabular}{lcccc}
\toprule
Learner (residual memory) & $\Delta$ intensity & $\Delta$ informativeness & Rival $\Delta\beta_1$ & Deviator gain \\
\midrule
Tabular Q, $\gamma=0.95$ & $0.86 \pm 0.01$ & $0.93 \pm 0.01$ & $-0.000 \pm 0.003$ & $0.053 \pm 0.003$ \\
Tabular Q, $\gamma=0$ & $1.24 \pm 0.00$ & $1.46 \pm 0.01$ & $0.002 \pm 0.003$ & $0.124 \pm 0.004$ \\
DQN, $\gamma=0.95$ & $-0.36 \pm 0.07$ & $-0.17 \pm 0.05$ & $-0.000 \pm 0.004$ & $0.026 \pm 0.008$ \\
DQN, $\gamma=0$ & $-0.17 \pm 0.03$ & $-0.09 \pm 0.02$ & $-0.001 \pm 0.005$ & $0.014 \pm 0.007$ \\
PPO, $\gamma=0.95$ & $-0.02 \pm 0.09$ & $0.06 \pm 0.07$ & $0.000 \pm 0.000$ & $0.028 \pm 0.008$ \\
PPO, $\gamma=0$ & $-0.28 \pm 0.05$ & $-0.18 \pm 0.04$ & $0.000 \pm 0.000$ & $0.012 \pm 0.006$ \\
\bottomrule
\end{tabular}}
  \caption{Learning algorithms in the standard Kyle market ($\xi=0$, $I=2$,
  residual memory). Tabular Q: 600 markets, $6\times10^6$ periods; DQN and
  PPO: 50 markets, $4\times10^5$ periods. 95\% CIs.}
  \label{tab:deep}
\end{table}
